\section{Discusión y conclusiones}
\label{sec:discusion}

Esta última sección junta lo que se puede afirmar y lo que no. Está escrita para que el lector pueda decidir con qué está de acuerdo, así que insiste más en los puntos donde se podría discrepar que en los resultados.

\subsection{Qué se ha establecido}

El resultado analítico es que, sobre un fondo FLRW plano con un inflatón y bajo la hipótesis de que la materia no aporta tensión anisótropa TT, la ecuación de los modos tensoriales en UG, incluso con difusión, es la de RG (ecuación \eqref{eq:modok}), y que su acción cuadrática lo es también (ecuación \eqref{eq:S2TT}), esta última como hipótesis explícita de la realización efectiva (la ecuación de propagación no fija por sí sola la normalización cuántica), y que la difusión $Q$ y la constante de integración $\Lambda_0$ entran en el espectro tensorial únicamente a través de la función $H(t)$ del fondo. La argumentación distingue la ecuación lineal y una expansión formal condicionada, una a nivel de las ecuaciones de campo (el multiplicador es un escalar y por lo tanto pura traza, y el vínculo se cumple sin condiciones sobre un tensor sin traza a primer orden) y otra a nivel de la acción (el término del vínculo a segundo orden, $+\Lb a^3h_{ij}^2/4\kap$, cancela exactamente el término de masa que dejarían los otros dos aportes). Ambas coinciden con lo que la literatura afirma sin derivarlo del todo \cite{basak16,cho15,fabris22,barvinsky19}. De lo derivado a mano, la ecuación de los tensores a primer orden se verificó con álgebra simbólica (sección~\ref{sec:simbTT}); la acción a segundo orden y la cancelación del vínculo no, y hay un coeficiente (el del término de gradiente de la acción) que se fijó por consistencia y no por derivación.

El resultado numérico es que, para un inflatón y una difusión $Q(N)=Q_i\ee^{-\gamma N}$ con $Q_i/V_i=0.1$ y $\gamma=0.1$, y con las mismas condiciones iniciales en ambos modelos, la diferencia entre el espectro tensorial de UG y el de RG depende del potencial y del criterio de comparación (secciones~\ref{sec:resultados} y~\ref{sec:potenciales}). Con el potencial cuadrático, a igual $k$ comóvil UG es entre un 8.6\% y un 65.6\% menor que RG, pero a igual número de e-folds antes del final de cada inflación es un 50\% \emph{mayor}: el signo no es una propiedad de UG. Con el potencial de Starobinsky la reducción es $9.6$--$9.8\%$ a igual $N_*$ y $5.3$--$20.4\%$ a igual $k$, y las diferencias de $n_T$ son pequeñas (no se evaluó su detectabilidad instrumental). Con el modelo de radiación con difusión Q \cite{leon22,piccirilli23} (radiación más $Q$, sección~\ref{sec:leon}) no hay diferencia tensorial a igual $H(N)$: $P_T$ es el de RG con ese $H(N)$. Las diferencias de las secciones~\ref{sec:resultados} y~\ref{sec:potenciales} no son artefactos numéricos y se explican por el fondo: el cociente numérico coincide con la fórmula estándar de RG evaluada con el $H$ y el $\epsilon_1$ de UG con máximo $2.80\times10^{-5}$ en los modos juzgados del cuadrático (muestra de C3.4; se excluyen los puntos que no satisfacen su criterio slow-roll). El código reproduce además, con una precisión muy por debajo de la diferencia UG--RG, el límite estándar de RG (de Sitter exacto, desarrollo de slow-roll a segundo orden, las predicciones conocidas de $\phi^2$).

\subsection{De qué decisiones depende la magnitud}
\label{sec:decisiones-dep}

Lo anterior se dice para \emph{estos} parámetros, y es importante entender cuánto de la magnitud de la diferencia es una propiedad de UG y cuánto es una propiedad de la parametrización elegida. La respuesta es que casi toda la magnitud es de lo segundo.

La decisión más fuerte es la de $\Lambda_0$. Pedir las mismas condiciones iniciales ($\phi_i$, $\dot\phi_i$, $H_i$) en RG y en UG fija $\Lambda_0=-Q_i$, negativa y comparable a la densidad de energía total al inicio ($Q_i$ es el 10\% de $V(\phi_i)$). Como mostró la figura~\ref{fig:fondo}, eso hace que $3H^2=\rho_\phi-Q_i+Q(N)$ pierda energía respecto de RG durante toda la inflación y que, al final, esa constante negativa sea comparable a $\rho_\phi$: la inflación de UG (con el potencial cuadrático) termina en $\epsilon_1=1$, aunque el parámetro del potencial desnudo sea pequeño, y $H$ cae mucho más rápido que en RG. En Starobinsky el presupuesto energético al final es distinto, pero no se aisló la causa variando la constante. Una elección alternativa igualmente razonable, $\Lambda_0=0$ con los mismos $(\phi_i,\dot\phi_i)$, daría $H_i^{\rm UG}=\sqrt{(\rho_i+Q_i)/3}>H_i^{\rm RG}$ y un $H$ mayor en UG mientras $Q$ es apreciable; a partir de \eqref{eq:H2UG} cabe esperar que el signo de la diferencia en $P_T$ sea el opuesto en la primera parte del rango de $k$. No se calculó, y no se afirma. Lo que sí se calculó (sección~\ref{sec:resultadosN}) es que, incluso con $\Lambda_0=-Q_i$, el signo cambia según se compare a igual $k$ o a igual $N_*$: la magnitud y el signo de la diferencia dependen de decisiones que no fija la teoría.

Otras decisiones actúan más suavemente. La forma de $Q(N)$ y los valores de $Q_i/V_i$ y $\gamma$ son elecciones fenomenológicas; un barrido exploratorio preliminar en el notebook (no verificado y no registrado como resultado) sugiere que el cociente UG/RG varía fuertemente con $Q_i/V_i$. La velocidad inicial de slow-roll de RG produce el salto inicial de $\epsilon_1$ de UG, que es una consecuencia de imponer las mismas condiciones iniciales y no una propiedad de la teoría. La duración de la inflación de UG ($100$ e-folds) se fijó por la literatura, y hace que RG dure más ($161.8$ con el potencial cuadrático y $386.7$ con el de Starobinsky). Y el mismo $k$ no corresponde a la misma escala observacional en los dos modelos; por eso la comparación principal es a igual $N_*$ (sección~\ref{sec:igualN}). Es una convención ilustrativa, no una identificación de la misma escala observacional: identificar el mismo $k_*$ físico exige la historia posterior (el recalentamiento) y la normalización de $a$, que no se modelan.

Lo que sí es robusto, siempre que la derivación sea correcta, es el \emph{mecanismo}: $P_T$ depende de $Q$ solo a través de $H(t)$, de modo que cualquier efecto de la difusión sobre el espectro tensorial se lee en el fondo. La cantidad de efecto es, en cambio, una predicción del modelo de $Q$ que se elija, y no de UG como tal.

\subsection{Supuestos que sostienen el resultado analítico}
\label{sec:supuestos}

Hay varios puntos en los que el lector podría legítimamente discrepar, y conviene enumerarlos.

El primero y más importante es la hipótesis sobre la materia. La sección~\ref{sec:acoplamiento} la precisa para el caso de un solo campo: si $Q$ intercambia energía con el inflatón y es covariante, para un escalar canónico y una difusión local especificada como $Q(\phi)$, las ecuaciones de UG pueden reescribirse como las de RG con $V_{\rm eff}=V+Q(\phi)+\Lambda_0$ (sobre la rama reconstruida, si se parte de $Q(N)$); y si además se impone $\delta Q=0$ en el marco especificado y la perturbación del inflatón es inhomogénea, resulta $\dot Q=0$. Ninguna de las dos afirmaciones es una equivalencia o una prohibición general (sección~\ref{sec:acoplamiento}). La conclusión de que la acción cuadrática de los tensores es la de RG usa que el tensor energía-impulso no tiene parte TT a primer orden, \eqref{eq:TTmateria}, y que el aporte de la materia a la acción a segundo orden es el estándar, \eqref{eq:Sm2}. La ausencia de parte TT a primer orden no justifica por sí sola el aporte de materia a segundo orden con intercambio de energía. Pero, como se dijo en el recuadro de la sección~\ref{sec:difusionQ}, la difusión $Q$ variable exige que la materia \emph{no} sea la de una acción covariante estándar, y no se sabe qué modificaciones concretas de la materia producen esa no conservación. Si el mecanismo microscópico detrás de $Q$ (por ejemplo, los efectos de \cite{josset17,perez19}) generara una tensión anisotrópica, o modificara la acción de la materia a segundo orden en $h_{ij}$, los tensores podrían verse afectados de un modo que este análisis no captura. Es un supuesto compartido con las demás referencias que tratan tensores con difusión \cite{fabris22}.

El segundo es el tratamiento del vínculo a segundo orden. La derivación fija $g_{00}=-1$ y deja que el multiplicador absorba la violación de segundo orden del volumen. Una alternativa es perturbar también el lapse a segundo orden, $N=1+\tfrac14h_{ij}^2$, de modo que el vínculo se cumpla exactamente; la diferencia entre ambas formulaciones multiplica la ecuación de fondo de la componente temporal, y cabe esperar que la acción en capa sea la misma, pero ese argumento no se desarrolló en el trabajo y sería bueno comprobarlo con álgebra simbólica. La observación de la sección~\ref{sec:accion2} sobre la parametrización exponencial cubre una elección relacionada, no esta.

El tercero es que el coeficiente del término de gradiente de \eqref{eq:S2TT}, $-\tfrac14a(\partial_kh_{ij})^2$, no se derivó sino que se tomó como estándar y se contrastó por consistencia con la ecuación de campo. Y el cuarto es más general: la derivación es a mano, y algunos pasos estándar (por ejemplo, la forma ADM de Einstein--Hilbert) no se contrastaron con el texto de un libro. La verificación simbólica cubre solo la ecuación de los tensores a primer orden y algunas identidades del fondo y del acoplamiento de $Q$ (sección~\ref{sec:acoplamiento}).

Hay además cosas que no se consideran. Se trabaja a nivel lineal en el fondo homogéneo, de modo que no se estudian las ondas gravitacionales inducidas a segundo orden por las perturbaciones escalares (en UG el sector escalar tiene un gauge restringido y $\delta Q$ podría contribuir), ni posibles correcciones al estado inicial de Bunch--Davies debidas a los mecanismos que generan $Q$.

\subsection{Lo que queda abierto}
\label{sec:abiertos}

La ausencia principal es el sector escalar, y con ella la relación tensor-escalar $r$, que es la cantidad observacionalmente relevante. Hay un problema específico en la literatura que este trabajo no resuelve pero que conviene describir, porque motivó el proyecto. El artículo de León \cite{leon22} da un espectro escalar (su Ec.~66) con un factor $3^{3/2}$ en el numerador, $P_\mathcal{R}=3^{3/2}H^2/(8\pi^2\MP^2\epsilon_1)$; el de Piccirilli y León \cite{piccirilli23}, que dice basarse en aquél, tiene la misma expresión sin ese factor (su Ec.~41), y luego adopta $r=16\,\epsilon_1$ (su Ec.~42) citando un trabajo que no contiene sector tensorial. Con el $P_\mathcal{R}$ de \cite{leon22} y el $P_T$ estándar, aritméticamente resultaría $r=16\epsilon_1/3^{3/2}\approx3.1\,\epsilon_1$, no $16\epsilon_1$. Esa observación, y una hipótesis sobre su origen (que $3^{3/2}$ podría venir de una velocidad del sonido $c_s^2=1/3$, aunque la expresión que recuerdo para un fluido con velocidad del sonido $c_s$, asociada a \cite{garriga99}, tiene $c_s^{-1}$, esto es $3^{1/2}$ y no $3^{3/2}$, y no la verifiqué contra el texto de ese artículo), son conjeturas mías que no se verificaron. Lo que este trabajo aporta a ese problema es la mitad tensorial: $P_T$ es el de RG evaluado con el $H$ de UG. La otra mitad, $P_\mathcal{R}$ con difusión, es una tarea pendiente.

Para el modelo con inflatón y $Q(\phi)$ el sector escalar tampoco está derivado: si la equivalencia con RG con $V_{\rm eff}$ fuese correcta a nivel de perturbaciones valdrían las fórmulas estándar, pero esa equivalencia se argumentó a mano y no se verificó (sección~\ref{sec:unescalar}). Y queda abierto el sector escalar en el marco $\delta Q=0$ con radiación, que es el de \cite{leon22,piccirilli23}; con el fondo de esa literatura, $A_s$ y $r$ dependen de una fórmula escalar no derivada, y el factor $3^{3/2}$ cambia $A_s$ por un factor cinco (sección~\ref{sec:leonescalar}). Otros puntos abiertos, en orden de prioridad sugerida: el escenario 2 del modelo de radiación con difusión Q \cite{leon22,piccirilli23}, que necesita $N_f\simeq370$ y una integración de precisión extendida (sección~\ref{sec:leon}), y una normalización del mapeo del escenario 3 que no acumule el error del final de la inflación; un barrido controlado de $Q_i$, $\gamma$ y $\phi_i$, y de la prescripción de $\Lambda_0$; la relación entre $k$ y la escala observacional, que requiere modelar el recalentamiento; la verificación simbólica de la acción a segundo orden y el contraste de los pasos estándar con los textos de referencia; y, por requisito del curso, el registro formal de los números y afirmaciones del trabajo con su procedencia y una segunda lectura independiente, hecha por alguien sin la historia de este trabajo.

\subsection{Conclusiones}

El trabajo intentó cerrar un hueco específico de la literatura sobre gravedad unimodular con difusión: partiendo de la acción con multiplicador de Lagrange, ¿cambia la ecuación de las ondas gravitacionales primordiales? La respuesta que se obtiene, con los supuestos explicitados arriba, es que no cambia: el vínculo unimodular no toca a los modos transversos y sin traza a primer orden, y a segundo orden su contribución a la acción se cancela exactamente con las demás, de modo que la difusión $Q$ solo llega al espectro tensorial a través del parámetro de Hubble. Los controles numéricos son consistentes con esa conclusión, aunque no la prueban de manera independiente.

Lo que sí se puede afirmar, además, es de naturaleza más práctica. El espectro tensorial en UG con difusión se calcula como en RG con el fondo de UG, y toda la información sobre $Q$ que llega al espectro está en $H(t)$; por lo tanto, medir $P_T$ (o acotar $r$) restringe la difusión $Q$ solo a través de su efecto sobre el fondo. Y cuánto restringe depende, como se vio, de decisiones de modelado (empezando por $\Lambda_0$, el potencial y el criterio de comparación) que la teoría no fija; en particular, para un escalar canónico con $Q(\phi)$ el efecto se reduciría a usar otro potencial, siempre que la equivalencia perturbativa se confirme (sección~\ref{sec:unescalar}). El mensaje central es entonces: a igual fondo, y bajo las hipótesis adoptadas, la propagación tensorial no distingue UG de RG; las diferencias de los ejemplos provienen del fondo y del criterio de comparación; el acoplamiento escalar y las fórmulas alternativas de $A_s$ son discusión condicionada. Quien lea este informe con ojo crítico debería decidir, en primer lugar, si le parecen aceptables los supuestos sobre la materia de la sección~\ref{sec:supuestos}; en segundo lugar, si la decisión de fijar el mismo $H_i$ en RG y en UG y de comparar a igual $N_*$ es la comparación que le interesa; y, en tercer lugar, si el marco de trabajo es el de un inflatón con $Q=Q(\phi)$ o el de $\delta Q=0$ con radiación, y si el siguiente paso es el sector escalar en ese marco o el barrido de parámetros.
